\documentclass[10pt, letterpaper]{article}

% --- PACKAGES AND SETUP ---
\usepackage[utf8]{inputenc}
\usepackage[T1]{fontenc}
\usepackage[english]{babel}
\usepackage[margin=1in]{geometry}
\usepackage{parskip}
\usepackage{amsmath}
\usepackage{amssymb}
\usepackage{soul}
\usepackage{csquotes}
\usepackage{ragged2e}
\usepackage[none]{hyphenat}
\usepackage{chngcntr}
\usepackage{enumitem}
\usepackage{caption}

% --- BIBLIOGRAPHY SETUP (biblatex) ---
\usepackage[
    backend=biber,
    style=authoryear,
    sorting=nyt,
    maxcitenames=2,
    mincitenames=1 
]{biblatex}

\addbibresource{grass.bib}
\counterwithin{subsection}{section}

% --- HYPERLINKS AND PDF METADATA ---
\usepackage{hyperref}
\hypersetup{
    colorlinks=true,
    linkcolor=blue,
    citecolor=blue,
    urlcolor=blue,
    pdftitle={The Aggregate Opportunity Cost of Recreational Land Use},
    pdfauthor={Michael Manolakis},
}

% --- DOCUMENT START ---
\begin{document}

\begin{center}
    {\huge The Aggregate Opportunity Cost of Recreational Land Use \\}
    {\large A Counterfactual Valuation of U.S. Golf Courses \\}
    {\normalsize Michael Manolakis \textbf{\textemdash} April 1, 2026 \\}
\end{center}

\raggedright

\vspace{1cm}

\begin{abstract}
    \textit{[Placeholder: 228-word summary detailing the macroeconomic inefficiency of golf courses, the HBU methodology, the MICE imputation of 16,297 facilities, the \$244.9B pooled result, and the institutional frictions of zoning and covenants.]}
\end{abstract}

\vspace{1cm}

\section{Introduction}

\subsection{Macroeconomic Context and Spatial Misallocation}
\begin{itemize}
    \item \textbf{The Problem:} Establish the inelasticity of urban land supply and the resulting "zoning tax" that inflates residential land values.
    \item \textbf{The Inefficiency:} Introduce the preservation of low-density recreational land (golf courses) in these constrained markets as a macroeconomic distortion.
\end{itemize}

\subsection{The Golf Course Typology and Amenity Depreciation}
\begin{itemize}
    \item \textbf{Joint Production:} Explain how golf courses historically functioned as surrogate public goods to cross-subsidize residential real estate premiums.
    \item \textbf{Shifting Preferences:} Note the structural decline in golf demand and the shift in consumer preference toward natural open spaces, leaving massive acreage locked in a depreciating use.
\end{itemize}

\subsection{Methodological Framework: Highest and Best Use (HBU)}
\begin{itemize}
    \item \textbf{The Counterfactual:} Introduce the appraisal doctrine of HBU and define the core opportunity cost equation: $V_{OC} = \max(V_{Res}, V_{Ag}) - V_{Current}$.
    \item \textbf{The Hypothetical Condition:} Explicitly state the assumption of a frictionless legal environment (dissolving zoning and private covenants) to isolate pure economic value.
\end{itemize}

\subsection{Research Question and Roadmap}
\begin{itemize}
    \item \textbf{Core Question:} What is the aggregate financial opportunity cost of U.S. golf courses when evaluated against their HBU counterfactuals?
    \item \textbf{Roadmap:} Briefly outline the progression of the thesis through literature, theory, data (MICE/RUCC), results, and policy implications.
\end{itemize}

\newpage

\section{Literature Review}
\begin{itemize}
    \item \textbf{Urban Land Valuation:} Discuss Davis et al. and Glaeser regarding the "zoning tax" and the extreme spatial concentration of residential land value.
    \item \textbf{Microeconomics of Amenities:} Discuss Rosen, Limehouse, and Cho regarding hedonic pricing, joint production, and the depreciating amenity value of golf courses.
    \item \textbf{Institutional Frictions:} Discuss Coase, Ellickson, and Korngold regarding the failure of market bargaining due to restrictive HOA covenants and exclusionary zoning.
\end{itemize}

\section{Theoretical Framework}
\begin{itemize}
    \item \textbf{Neoclassical Land Allocation:} Define spatial equilibrium and the mathematical maximization of residual land value.
    \item \textbf{The Opportunity Cost Function:} Deconstruct the variables $V_{Res}$, $V_{Ag}$, and $V_{Current}$.
    \item \textbf{Market Failure:} Model the institutional frictions $F(Z, C)$ (Zoning and Covenants) that prevent $V_{OC}$ from being realized, justifying the Hypothetical Condition.
\end{itemize}

\section{Data and Methodology}
\begin{itemize}
    \item \textbf{Dataset Scope:} Detail the 16,297 observations, Tigris spatial boundaries, FHFA residential data, and USDA agricultural data.
    \item \textbf{Hybrid Valuation Algorithm:} Explain the use of USDA RUCC codes to bifurcate the data into Urban ($V_{Res}$) and Rural ($V_{Ag}$) counterfactuals to control for elasticity.
    \item \textbf{MICE Imputation:} Detail the Predictive Mean Matching (PMM) protocol used to resolve the 38\% missing acreage data across $m=5$ datasets.
    \item \textbf{Parameter Pooling:} Explain the use of Rubin's Rules to pool the final aggregate sums and regression coefficients to maintain statistical validity.
\end{itemize}

\section{Empirical Results}
\begin{itemize}
    \item \textbf{Aggregate Valuation:} Present the pooled \$244.9 billion opportunity cost and its 95\% confidence interval.
    \item \textbf{Spatial Heterogeneity:} Present the Urban vs. Rural bifurcation table, highlighting that 98.5\% of the value is concentrated in urban areas.
    \item \textbf{Econometric Determinants:} Present the pooled logarithmic regression output, interpreting the coefficients for holes, course type (Public/Private), and urban classification.
\end{itemize}

\newpage

\section{Discussion and Policy Implications}
\begin{itemize}
    \item \textbf{Macroeconomic Impact:} Translate the \$244.9B figure into the real-world cost of foregone housing supply and misallocated labor.
    \item \textbf{The Legal Straitjacket:} Discuss the NIMBY resistance and the "stale covenants" that legally paralyze the land, making Coasian bargaining impossible.
    \item \textbf{Alternative Counterfactuals:} Introduce Weinand's research on utility-scale renewable energy (solar/wind) as a highly synergistic, non-residential HBU.
    \item \textbf{Methodological Limitations:} Acknowledge the "bulk discount" (plattage) limitation of applying subdivided residential lot prices to massive 150-acre parcels.
\end{itemize}

\section{Conclusion}
\begin{itemize}
    \item \textbf{Summary:} Restate the \$244.9B finding, the MICE/RUCC methodology, and the extreme urban concentration of the opportunity cost.
    \item \textbf{Final Thought:} Conclude that unlocking this value requires aggressive legal reform to overcome the institutional frictions of zoning and private covenants.
\end{itemize}

\newpage

\printbibliography[title={References}]

\end{document}